\documentclass[../../main.tex]{subfiles}
\begin{document}

{\bf[解]}

\begin{figure}[htb]
  \centering
  \begin{tikzpicture}[scale=2, >=stealth, every node/.style={font=\small}]
    % t=0.8，a=t+1/(2t^2)，b=1/t-1/(4t^4)の例．
    \pgfmathsetmacro{\contactT}{0.8}
    \pgfmathsetmacro{\shiftA}{\contactT+1/(2*\contactT^2)}
    \pgfmathsetmacro{\shiftB}{1/\contactT-1/(4*\contactT^4)}
    \coordinate (T) at (\contactT,{1/\contactT});
    \draw[->,thick] (-0.2,0) -- (3.4,0) node[right] {$x$};
    \draw[->,thick] (0,-0.3) -- (0,3.4) node[above] {$y$};
    \node[below left] at (0,0) {$O$};
    \begin{scope}
      \clip (0.04,-0.2) rectangle (3.15,3.1);
      \draw[very thick,blue!85!black,domain=0.05:3.15,samples=101]
      plot (\x,{(\x-\shiftA)^2+\shiftB});
      \draw[very thick,teal!85!black,domain=0.25:3.15,samples=121]
      plot (\x,{1/\x});
      \draw[thick,red!85!black,domain=0.05:1.72]
      plot (\x,{2/\contactT-\x/(\contactT^2)});
    \end{scope}
    \draw[dashed,gray!60] (\contactT,0) -- (T) -- (0,{1/\contactT});
    \fill (T) circle (0.8pt);
    \node[above right] at (T) {$\Bigl(t,\frac{1}{t}\Bigr)$};
    \node[below] at (\contactT,0) {$t$};
    \node[left] at (0,{1/\contactT}) {$\frac{1}{t}$};
    \node[right,blue!85!black] at (2.75,2.20) {$C$};
    \node[above,teal!85!black] at (2.65,0.40) {$C'$};
    \node[right,red!85!black] at (1.38,0.22) {共通接線};
  \end{tikzpicture}
  \caption{放物線$C$と曲線$C'$が$x=t$で接し，共通接線をもつ様子}
  \label{fig:utokyo1987q2_common_tangent}
\end{figure}

\paragraph{$(a,b)$の範囲の図示}

$y=x^2$を$(a,b)$だけ並行移動して，$C$の方程式は
\begin{align}
  C:y=(x-a)^2+b \label{eq:1}
\end{align}
である．
また
\begin{align}
  C': y = \dfrac{1}{x} \label{eq:2}
\end{align}
である．
\cref{eq:1,eq:2}の$x=t\,(>0)$での接線が一致する時
$\left(\cfrac{1}{x}\right)'=\cfrac{-1}{x^2}$だから条件は，
\begin{align}
  \begin{cases}
    t>0 \\
    (t-a)^2+b=\cfrac{1}{t} \\
    2(t-a)=\cfrac{-1}{t^2}
  \end{cases}
  \Longleftrightarrow
  \begin{cases}
    t>0 \\
    b=\cfrac{1}{t}-\left(\cfrac{-1}{2t^2}\right)^2 \\
    a=t+\cfrac{1}{2t^2}
  \end{cases}
  \Longleftrightarrow
  \begin{cases}
    t>0 \\
    b=\cfrac{1}{t}-\cfrac{1}{4t^4} \\
    a=t+\cfrac{1}{2t^2} \label{eq:3}
  \end{cases}
\end{align}
なる$t$が存在することである．
従って，$(a,b)$が\cref{eq:3}で表されるパラメータ曲線上にあればよい．
$a,b$を$t$で一階微分すると
\begin{align}
  \begin{cases}
    a'=1-\dfrac{1}{t^3}=\dfrac{t^3-1}{t^3} \\
    b'=\dfrac{-1}{t^2}+\dfrac{1}{t^5}=\dfrac{1-t^3}{t^5}
  \end{cases}
\end{align}
だから下表を得る．

\begin{table}[htb]
  \centering
  \caption{$a(t)$，$b(t)$ の増減表}
  \label{tab:utokyo1987q2_increase_decrease}
  \renewcommand{\arraystretch}{1.4}
  \begin{tabular}{c|c|c|c|c|c}
    $t$ & $(0)$ & $\dots$ & $1$ & $\dots$ & $(+\infty)$ \\ \hline
    $a'$ & & $-$ & $0$ & $+$ & \\ \hline
    $b'$ & & $+$ & $0$ & $-$ & \\ \hline
    $(a,b)$ & & $\nwarrow$ & $\Bigl(\dfrac{3}{2},\dfrac{3}{4}\Bigr)$ & $\searrow$ &
  \end{tabular}
\end{table}

また，\cref{eq:3}から極限値は以下のようになる．
\begin{align}
  t\to+0\text{のとき}&
  \begin{cases}
    a\to+\infty \\
    b\to-\infty
  \end{cases}\\
  t\to+\infty\text{のとき}&
  \begin{cases}
    a\to+\infty \\
    b\to+0
  \end{cases}
\end{align}

従って，$(a,b)$のグラフの概形は以下のようになる．

\begin{figure}[htb]
  \centering
  \begin{tikzpicture}[scale=1.25, >=stealth, every node/.style={font=\small}]
    \draw[->] (-0.25,0) -- (5.2,0) node[right] {$a$};
    \draw[->] (0,-2.45) -- (0,1.35) node[above] {$b$};
    \node[below left] at (0,0) {$O$};
    \begin{scope}
      \clip (0.1,-2.3) rectangle (5,1.15);
      \draw[very thick, blue!85!black, domain=0.45:1, samples=121, variable=\t]
      plot ({\t+1/(2*\t^2)},{1/\t-1/(4*\t^4)});
      \draw[very thick, blue!85!black, domain=1:5.1, samples=121, variable=\t]
      plot ({\t+1/(2*\t^2)},{1/\t-1/(4*\t^4)});
    \end{scope}
    \draw[dashed, gray!60] (1.5,0) -- (1.5,0.75) -- (0,0.75);
    \fill[blue!85!black] (1.5,0.75) circle (0.8pt);
    \node[below] at (1.5,0) {$\frac{3}{2}$};
    \node[left] at (0,0.75) {$\frac{3}{4}$};
    \node[above right] at (1.5,0.75) {$\Bigl(\frac{3}{2},\frac{3}{4}\Bigr)$};
    \node[above, blue!85!black] at (3.5,0.42) {$t>1$};
    \node[right, blue!85!black] at (2.65,-1.7) {$0<t<1$};
  \end{tikzpicture}
  \caption{$a=t+\dfrac{1}{2t^2}$，$b=\dfrac{1}{t}-\dfrac{1}{4t^4}$（$t>0$）で表される$(a,b)$の軌跡}
  \label{fig:utokyo1987q2_locus}
\end{figure}

\paragraph{接点以外に共有点を持たないような$(a,b)$の値}

次に，接点以外に交点を持たない時を考えよう．これは\cref{eq:1,eq:2}を連立した方程式
\begin{align}
  (x-a)^2+b=\cfrac{1}{x} \quad (x>0) \label{eq:4}
\end{align}
を満たす$x$が唯一つ$x=t$のみであればよい．

\cref{eq:3}を代入して
\begin{align}
  &\left(x-\left(t+\cfrac{1}{2t^2}\right)\right)^2+\left(\cfrac{1}{t}-\cfrac{1}{4t^4}\right)=\cfrac{1}{x} \\
  \Longleftrightarrow%
  &x^2-2\left(t+\cfrac{1}{2t^2}\right)x+\left(t+\cfrac{1}{2t^2}\right)^2+\left(\cfrac{1}{t}-\cfrac{1}{4t^4}\right)=\cfrac{1}{x} \\
  \Longleftrightarrow%
  &x^3-2\left(t+\cfrac{1}{2t^2}\right)x^2+\left(t^2+\cfrac{2}{t}\right)x-1=0 \\
  \Longleftrightarrow%
  &(x-t)^2\left(x-\cfrac{1}{t^2}\right)=0
\end{align}
だから，この解が$x=t$のみである条件は$t>0$も考慮して
\begin{align}
  t=\cfrac{1}{t^2}\Longleftrightarrow t=1
\end{align}
である．
この時\cref{eq:3}から，$(a,b)=\left(\dfrac{3}{2},\dfrac{3}{4}\right)$となる．$\cdots$(答)
\newpage
\end{document}
